% =========================================================
% SECCIÓN 3.5 - CINÉTICA DE ASIMILACIÓN Y MODULACIÓN AMBIENTAL
% =========================================================

Las ecuaciones de balance del sistema de EDO formuladas en la Sección \ref{sec:cap3_formulacion_ode} requieren la especificación explícita de la tasa de ingesta $\tilde{\imath}(S,B,T)$ y de la tasa de mantenimiento $\dot{M}(B,T)$. Esta sección cierra esas formas funcionales mediante una cinética de ingesta tipo Holling II —que captura la saturación por disponibilidad de sustrato— y un factor de dependencia térmica basado en el modelo de Sharpe--Schoolfield —que describe la modulación de la tasa metabólica por la temperatura del sustrato. Ambas funciones se parametrizan con semillas tomadas de la literatura y dejan variables libres para su ajuste local en el Capítulo \ref{chap:hibrido}.

% ---------------------------------------------------------
\subsection{Funcional respuesta tipo Holling II}
\label{subsec:cap3_holling}
% ---------------------------------------------------------

La ingesta específica por larva, $\tilde{\imath}$ (g C/larva/min), se modela como el producto de tres factores multiplicativos: una tasa máxima de procesamiento, una función de disponibilidad de sustrato y un factor de modulación térmica:
\begin{equation}
    \tilde{\imath}\bigl(S(t), B(t), T(t)\bigr) = a_{\max}\, B(t)^{2/3}\, f\bigl(T(t)\bigr)\, \frac{S(t)}{\tau_h + S(t)}.
    \label{eq:cap3_holling}
\end{equation}

El factor $B^{2/3}$ refleja la hipótesis de similitud isométrica: la superficie bucal y el área de contacto con el sustrato crecen proporcionalmente al cuadrado del diámetro corporal, mientras que la masa estructural crece con el cubo; por ende, la tasa de ingesta específica escala con la superficie relativa $B^{2/3}$ \parencite{eriksenDynamicModellingFeed2022}. Esta aproximación es consistente con los rangos de densidad larval del protocolo experimental (\SI{2.5}{}--\SI{5.0}{\gram\per\litre}), donde la competencia por el alimento no altera sustancialmente la geometría de contacto individual.

El parámetro $a_{\max}$ (g C$^{1/3}$/min) representa la tasa máxima de búsqueda y procesamiento de alimento por unidad de superficie metabólica a la temperatura de referencia. \textcite{eriksenDynamicModellingFeed2022} reportan un valor semilla de $a_{\max} \approx 0.18$ g C$^{1/3}$/día para larvas de BSF en sustrato de alta calidad, equivalente a $1.25 \times 10^{-4}$ g C$^{1/3}$/min tras conversión de unidades.

La fracción $S/(\tau_h + S)$ es la funcional respuesta tipo Holling II, donde $\tau_h$ (g C) es la constante de semisaturación del sustrato —esto es, la masa de carbono disponible a la cual la ingesta alcanza la mitad de su valor máximo. Este parámetro se deja como variable libre de calibración ($\tau_h \in [50, 500]$ g C) dado que depende fuertemente de la composición física del sustrato (partículas gruesas de pulpa de café versus cáscaras de fruta desmenuzadas) y de la humedad, que afecta la plasticidad del medio y la eficiencia de fragmentación mecánica por las larvas \parencite{bekkerImpactSubstrateMoisture2021}.

La asimilación neta por larva se obtiene multiplicando la ingesta por la eficiencia de asimilación intestinal $\kappa_A$:
\begin{equation}
    \tilde{a}\bigl(S(t), B(t), T(t)\bigr) = \kappa_A \cdot \tilde{\imath}\bigl(S(t), B(t), T(t)\bigr),
    \label{eq:cap3_asimilacion_neta}
\end{equation}
con $\kappa_A \in [0.50, 0.80]$ según la calidad del sustrato \parencite{guillaumeAsymptoticEstimatedDigestibility2023}. En el protocolo experimental del Capítulo \ref{chap:instrumentacion}, la mezcla de pulpa de café y cáscaras de frutas ($70:30$) se clasifica como sustrato de alta aptitud, por lo que se adopta la semilla $\kappa_A = 0.70$.

La ingesta total del reactor resulta entonces $\dot{I}_{\mathrm{tot}} = N_{\mathrm{larvas}} \cdot \tilde{\imath}$, que ingresa directamente en la ecuación \eqref{eq:cap3_ode_sustrato} del balance de sustrato. La no linealidad de la Holling II garantiza que, cuando $S \gg \tau_h$, la ingesta se satura en $a_{\max} B^{2/3} f(T)$, evitando predicciones no realistas de consumo infinito; conversamente, cuando $S \to 0$, la ingesta se anula y el sustrato se preserva no negativo.

% ---------------------------------------------------------
\subsection{Dependencia térmica: modelo de Sharpe--Schoolfield}
\label{subsec:cap3_sharpe_schoolfield}
% ---------------------------------------------------------

La temperatura del sustrato constituye la variable ambiental de mayor impacto sobre la cinética metabólica de \emph{H. illucens}. Estudios de campo y de laboratorio documentan que la tasa de asimilación y el crecimiento larval se incrementan de forma monótona desde los \SI{20}{\celsius} hasta un óptimo cercano a los \SI{30}{}--\SI{32}{\celsius}, para luego colapsar abruptamente por encima de los \SI{35}{\celsius} debido a la desnaturalización enzimática y al estrés térmico \parencite{nayakHermetiaIllucensDiptera2024,salamEffectDifferentEnvironmental2022}.

Para capturar este perfil asimétrico, se adopta el modelo de Sharpe--Schoolfield con inactivación de alta temperatura, expresado en escala Kelvin:
\begin{equation}
    f(T) = \frac{\exp\!\left(\dfrac{T_A}{T_{\mathrm{ref}}} - \dfrac{T_A}{T}\right)}{1 + \exp\!\left(\dfrac{T_A}{T_H} - \dfrac{T_A}{T}\right)},
    \label{eq:cap3_sharpe_schoolfield}
\end{equation}
donde:
\begin{itemize}
    \item $T$ es la temperatura absoluta del sustrato (K), derivada de la lectura del sensor SHT30 inmerso a \SI{5}{\centi\metre} de profundidad (Sección \ref{sec:cap2_sensores_ambientales});
    \item $T_{\mathrm{ref}} = 303.15$ K (\SI{30}{\celsius}) es la temperatura de referencia del protocolo experimental;
    \item $T_A$ (K) es la energía de activación de Arrhenius, que gobierna el ascenso exponencial de la tasa metabólica por debajo del óptimo;
    \item $T_H$ (K) es la temperatura de inactivación enzimática por desnaturalización, que controla el colapso del metabolismo por encima del óptimo.
\end{itemize}

El modelo \eqref{eq:cap3_sharpe_schoolfield} asume que la inactivación a baja temperatura es despreciable en el rango operativo del reactor (\SI{25}{}--\SI{35}{\celsius}), hipótesis validada para insectos tropicales donde el umbral inferior de actividad metabólica se sitúa por debajo de los \SI{20}{\celsius} \parencite{chiaThresholdTemperaturesThermal2018}. En el rango experimental, la función $f(T)$ presenta un máximo suave en $T \approx T_{\mathrm{opt}}$, definido implícitamente por la condición $\mathrm{d}f/\mathrm{d}T = 0$, que ocurre cuando los términos de activación e inactivación se balancean.

Los parámetros térmicos se tratan como variables libres de calibración con rangos semilla basados en la literatura entomológica:
\begin{equation}
    T_A \in [2000, 8000]\ \mathrm{K}, \qquad
    T_H \in [305, 315]\ \mathrm{K}\ (\SI{32}{}--\SI{42}{\celsius}), \qquad
    T_{\mathrm{opt}} \in [301, 308]\ \mathrm{K}\ (\SI{28}{}--\SI{35}{\celsius}).
    \label{eq:cap3_rangos_termicos}
\end{equation}

La temperatura óptima $T_{\mathrm{opt}}$ no es un parámetro independiente en la ecuación \eqref{eq:cap3_sharpe_schoolfield}, sino una propiedad emergente de $T_A$ y $T_H$; sin embargo, se incluye como variable de calibración auxiliar en el algoritmo de optimización del Capítulo \ref{chap:hibrido} para restringir la búsqueda a regiones biológicamente plausibles.

La tasa de mantenimiento específica, que aparece en la ecuación \eqref{eq:cap3_mantenimiento}, se modula térmicamente mediante el mismo factor $f(T)$:
\begin{equation}
    \dot{M}(t) = m \cdot B(t) \cdot f\bigl(T(t)\bigr),
    \label{eq:cap3_mantenimiento_termico}
\end{equation}
con $m$ la tasa de mantenimiento a la temperatura de referencia. Esta formulación garantiza que tanto la ingesta como el mantenimiento comparten la misma dependencia térmica, manteniendo la coherencia bioenergética del modelo DEB \parencite{eriksenDynamicModellingFeed2022}.

Respecto a la humedad relativa (HR), la literatura indica que valores por debajo del $50\%$ reducen la tasa de asimilación y aumentan la mortalidad, mientras que valores superiores al $80\%$ favorecen la formación de microambientes anaeróbicos \parencite{bekkerImpactSubstrateMoisture2021}. En el presente modelo, la HR no se incorpora como variable cinética explícita en las EDO, sino que se trata como una condición de contorno operativa: el protocolo del Capítulo \ref{chap:instrumentacion} la mantiene dentro del rango $60\%$--$80\%$, y las desviaciones transitorias se absorben en la calibración de $\tau_h$ y $a_{\max}$. Esta simplificación es consistente con el enfoque mecanicista del DEB, donde la humedad afecta indirectamente la disponibilidad de sustrato pero no la termodinámica intrínseca del metabolismo larval.

En síntesis, las ecuaciones \eqref{eq:cap3_holling} y \eqref{eq:cap3_sharpe_schoolfield} cierran las formas funcionales de los flujos metabólicos, transformando las EDO del modelo DEB en un sistema predictivo que responde a las variables ambientales medidas por los sensores IoT del Capítulo \ref{chap:instrumentacion}. La siguiente sección detalla los parámetros numéricos, sus unidades, rangos semilla y las restricciones físicas que garantizan la conservación de masa.
